
\maketitle

\begin{abstract}
Delay embedding parameter selection relies on diagnostics ---
false nearest neighbours, Jacobian-trace convergence, prediction
error --- that implicitly assume deterministic, noise-free
dynamics.  We show empirically that these criteria share a noise
threshold, above which they lose diagnostic value: for the
Lorenz-63 system, false nearest neighbour percentages at the
correct embedding dimension rise from $2\%$ (clean) to $62\%$
($5\%$ noise), while the Jacobian trace becomes uninformative.

We propose residual autocorrelation convergence --- the decrease
of lag-1 autocorrelation in local-linear-fit residuals as
embedding dimension increases --- as a noise-robust alternative.
This criterion correctly identifies sufficient embedding dimension
under $5$--$20\%$ observational noise where deterministic criteria
fail, and provides a U-shaped minimum for lag selection.  Results
are validated on the R\"{o}ssler system.

The computational primitive --- local linear regression in delay
coordinates --- is shared with S-maps~\cite{Sugihara1994} and
Jacobian-based Lyapunov estimation~\cite{SanoSawada1985}.  The
contribution is the application to embedding assessment and the
characterization of noise operating regimes.
\end{abstract}

\begin{quote}\small
\textbf{Keywords.}
delay embedding, embedding dimension, noise robustness, local
linear model, false nearest neighbours, residual diagnostics,
time series analysis.
\end{quote}

% ------------------------------------------------------------------
\section{Introduction}
\label{sec:intro}
% ------------------------------------------------------------------

Delay embedding~\cite{Takens1981, SauerYorkeCasdagli1991}
reconstructs the state space of a dynamical system from a scalar
time series $\{y_t\}$ via delay vectors
\begin{equation}\label{eq:delay}
  \mathbf{z}_t = (y_t,\, y_{t-L},\, \ldots,\, y_{t-(d-1)L})
  \in \mathbb{R}^d.
\end{equation}
The embedding dimension $d$ and lag $L$ must be chosen from data.
Standard criteria include false nearest
neighbours~\cite{KennelBrownAbarbanel1992} (FNN), mutual
information for lag selection~\cite{FraserSwinney1986},
prediction error minimisation~\cite{RagwitzKantz2002}, and
convergence of Lyapunov exponent
estimates~\cite{SanoSawada1985}.

These criteria assume the underlying dynamics are deterministic
and the observations are noise-free.  Under observational noise,
they degrade.  FNN in particular is known to be sensitive to
noise~\cite{RagwitzKantz2002}, but the degradation has not been
systematically characterised.

We make three contributions:
\begin{enumerate}
  \item We characterise the noise threshold at which
    Jacobian-based criteria (including FNN and Lyapunov
    convergence) lose diagnostic value
    (\S\ref{sec:threshold}).
  \item We propose residual autocorrelation convergence as a
    noise-robust embedding criterion for both dimension and lag
    selection (\S\ref{sec:residual}).
  \item We explain the threshold via the distinction between
    deterministic and stochastic modelling commitments
    (\S\ref{sec:discussion}).
\end{enumerate}

% ------------------------------------------------------------------
\section{Method}
\label{sec:method}
% ------------------------------------------------------------------

At each point $\mathbf{z}$ in the reconstructed space, we fit a
local linear model
\begin{equation}\label{eq:local-linear}
  \mathbf{z}_{t+1} \approx A(\mathbf{z})\,
    (\mathbf{z}_t - \mathbf{z}) + \mathbf{b}(\mathbf{z})
\end{equation}
by weighted least squares with a compactly supported tricube
kernel $w(u) = (1 - u^3)^3$ for $\|u\| \leq 1$.  Neighbours
within bandwidth $h$ are identified via KD-tree.  The Jacobian
$A(\mathbf{z}) \in \mathbb{R}^{d \times d}$ is the local
linearisation of the pushforward map.

This computation is shared with S-maps~\cite{Sugihara1994,
DeyleSugihara2011} (used for prediction and causality detection)
and Jacobian-based Lyapunov
estimation~\cite{SanoSawada1985, EckmannKamphorstRuelleEtAl1986}
(used to compute dynamical invariants).  We apply it to a
different question: embedding assessment.

From the local fit we extract two diagnostics:

\medskip
\noindent\textbf{Deterministic diagnostic:} the Jacobian trace
\begin{equation}\label{eq:div}
  \mathrm{div}\, F(\mathbf{z}) = \mathrm{tr}(A(\mathbf{z})),
\end{equation}
averaged over the attractor.  For a faithful embedding of a
deterministic system, $\langle\mathrm{tr}(A)\rangle$ converges
to the sum of local Lyapunov exponents as $d$ increases.
Instability of $\langle\mathrm{tr}(A)\rangle$ with $d$ indicates
insufficient embedding.

\medskip
\noindent\textbf{Residual diagnostic:} the lag-1 autocorrelation
of residuals
\begin{equation}\label{eq:residual}
  \varepsilon_t = \mathbf{z}_{t+1} -
    \bigl(A(\mathbf{z}_t)\,(\mathbf{z}_t - \mathbf{z}_{\mathrm{nn}})
    + \mathbf{b}(\mathbf{z}_{\mathrm{nn}})\bigr),
\end{equation}
where $\mathbf{z}_{\mathrm{nn}}$ is the nearest query point.  If
the embedding is sufficient and the local linear model captures the
deterministic dynamics, $\{\varepsilon_t\}$ should be white noise.
Remaining autocorrelation indicates deterministic structure not
captured by the embedding.

The criterion: \emph{the embedding dimension is sufficient when
the residual lag-1 autocorrelation stabilises near zero.}  For lag
selection: \emph{the optimal lag minimises the residual
autocorrelation at sufficient $d$.}

% ------------------------------------------------------------------
\section{Noise threshold for deterministic criteria}
\label{sec:threshold}
% ------------------------------------------------------------------

We generate a Lorenz-63 trajectory ($N = 50\,000$, $\Delta t =
0.01$) and observe the $x$-component corrupted by additive
Gaussian noise at amplitudes $0$--$20\%$ of the signal standard
deviation.  Delay vectors are formed with lag $L = 10$.

\subsection*{False nearest neighbours}

Table~\ref{tab:fnn} shows FNN percentages (threshold $R = 2$) at
the correct embedding dimension $d = 3$ as a function of noise
level.

\begin{table}[h]
\centering
\begin{tabular}{@{}lcccccc@{}}
\toprule
Noise (\%) & 0 & 1 & 2 & 5 & 10 & 20 \\
\midrule
FNN at $d = 3$ (\%) & 1.8 & 27.0 & 40.6 & 62.4 & 74.5 & 79.4 \\
\bottomrule
\end{tabular}
\caption{False nearest neighbour percentage at the correct embedding
dimension $d = 3$ for the Lorenz-63 system under varying
observational noise.  Above $\sim 2\%$ noise, FNN incorrectly
rejects the true dimension.}
\label{tab:fnn}
\end{table}

\noindent At $5\%$ noise, FNN reports $62\%$ false neighbours at
the correct dimension --- indistinguishable from an insufficient
embedding.  The criterion becomes useless.

\subsection*{Jacobian trace}

The mean $|\langle\mathrm{tr}(A)\rangle|$ across the attractor
shows a similar threshold.  For clean data, it increases from
$0.7$ ($d = 2$) to $47.7$ ($d = 6$), clearly indicating
dimension convergence.  At $5\%$ noise, the range collapses to
$0.02$--$0.68$: no diagnostic signal remains.

The transition is sharp: at $2\%$ noise, diagnostic signal
persists (range $0.7$--$7.2$); at $5\%$, it vanishes.

% ------------------------------------------------------------------
\section{Residual autocorrelation as embedding criterion}
\label{sec:residual}
% ------------------------------------------------------------------

\subsection*{Dimension selection}

Table~\ref{tab:acf} shows the lag-1 residual autocorrelation as a
function of embedding dimension, for both clean and $5\%$-noisy
Lorenz data.

\begin{table}[h]
\centering
\begin{tabular}{@{}lcccccc@{}}
\toprule
$d$ & 2 & 3 & 4 & 5 & 6 & 7 \\
\midrule
Clean: acf(1) & 0.38 & 0.24 & 0.28 & 0.20 & 0.11 & 0.11 \\
Noisy (5\%): acf(1) & 0.36 & 0.13 & 0.14 & 0.11 & 0.11 & 0.06 \\
\bottomrule
\end{tabular}
\caption{Lag-1 residual autocorrelation versus embedding dimension.
Both clean and noisy data show a decreasing trend, stabilising at
$d \approx 5$--$6$.  The criterion works under noise where FNN and
Jacobian trace do not.}
\label{tab:acf}
\end{table}

\noindent The residual autocorrelation decreases with $d$ and
stabilises, correctly indicating sufficient embedding.  This holds
across the full noise range tested ($0$--$20\%$; see
Figure~\ref{fig:phase}).

\subsection*{Lag selection}

At fixed $d = 3$, the residual autocorrelation as a function of
lag $L$ shows a U-shaped profile with a minimum near $L \approx
7$--$11$ (corresponding to $0.07$--$0.11$ orbital periods).  This
minimum persists under $5\%$ noise, providing a noise-robust lag
selection criterion where mutual information has limited
applicability.

\subsection*{Validation on R\"{o}ssler}

The same patterns hold for the R\"{o}ssler system ($a = 0.2$, $b =
0.2$, $c = 5.7$): residual autocorrelation decreases with $d$ for
both clean (0.14 $\to$ 0.03) and noisy data (0.09 $\to$ 0.01),
while the Jacobian trace diagnostic works only for clean data.

% ------------------------------------------------------------------
\section{Discussion}
\label{sec:discussion}
% ------------------------------------------------------------------

The shared noise vulnerability of FNN, Jacobian trace, and
Lyapunov convergence is not coincidental.  All three test for
properties of deterministic dynamics: FNN tests for injectivity of
the delay map; the Jacobian trace tests for well-defined local
expansion rates; Lyapunov convergence tests for stability of
dynamical invariants.  Observational noise destroys injectivity
(nearby points in delay space may be far apart in state space due
to noise), corrupts Jacobian estimates, and introduces spurious
Lyapunov exponents.

The residual criterion survives because it tests a different
property: whether the local linear model has captured all
\emph{predictable} structure, regardless of noise.  Residual
autocorrelation measures remaining deterministic signal in the
prediction errors --- a property that is well-defined even when
the dynamics are stochastic.  The embedding is sufficient when no
deterministic structure remains in the residuals.

This distinction corresponds to a difference in modelling
commitment.  Deterministic criteria implicitly assume $\mathbf{z}_{t+1}
= F(\mathbf{z}_t)$ (single-valued dynamics); noise violates this
assumption and the criteria respond to the violation rather than
to embedding insufficiency.  The residual criterion assumes
$\mathbf{z}_{t+1} = F(\mathbf{z}_t) + \sigma(\mathbf{z}_t)
\boldsymbol{\eta}_t$ (drift plus diffusion); noise is built into
the model, and the diagnostic tests only for remaining
deterministic structure.

The practical recommendation: for clean data ($< 2\%$ noise),
deterministic criteria (FNN, Jacobian convergence) provide sharp
diagnostics.  Above this threshold, use residual autocorrelation
convergence for dimension selection and the U-shaped minimum for
lag selection.

\subsection*{Relation to existing methods}

The local linear computation is shared with
S-maps~\cite{Sugihara1994, DeyleSugihara2011, CenciSugiharaSaavedra2019}
(prediction and causality) and Ragwitz--Kantz~\cite{RagwitzKantz2002}
(embedding selection via prediction error).  Ragwitz--Kantz
minimise out-of-sample prediction error over $(d, L)$, which is
prediction-optimised; our residual autocorrelation criterion is
sufficiency-based.  These can disagree: prediction skill may
plateau before residuals whiten (diagnostic underfitting) or
overfit in noisy regimes.

Cenci et al.~\cite{CenciSugiharaSaavedra2019} regularise S-maps
under noise; regularisation of the local Jacobian may extend the
operating regime of the deterministic diagnostic, but this is not
explored here.

\subsection*{Limitations}

The noise threshold ($\sim 2$--$5\%$) is characterised empirically
for two systems (Lorenz-63, R\"{o}ssler).  Its dependence on
attractor dimension, data length, and bandwidth is not
systematically explored.  The residual criterion requires
sufficient data to estimate local linear maps reliably; in very
high dimensions or with sparse data, both criteria may fail.

\bibliographystyle{plain}
\bibliography{references}
